\documentclass[11pt,a4paper]{article}
\usepackage[T1]{fontenc}
\usepackage[utf8]{inputenc}
\usepackage{amsmath,amssymb,amsthm}
\usepackage[margin=1in]{geometry}
\usepackage{booktabs}
\usepackage[colorlinks=true,linkcolor=blue,urlcolor=blue]{hyperref}
\theoremstyle{plain}
\newtheorem{theorem}{Theorem}
\newtheorem{lemma}[theorem]{Lemma}
\newtheorem{proposition}[theorem]{Proposition}
\newtheorem{corollary}[theorem]{Corollary}
\theoremstyle{definition}
\newtheorem{remark}[theorem]{Remark}

\title{The empirical recurrence for OEIS A196074 is false at every index its b-file can
test, and the correct one, of order 43}
\author{Adrian Perez Fontelles\\ \small Independent researcher}
\date{22 September 2026}
\begin{document}
\maketitle

\begin{abstract}
OEIS A196074 carries an empirical recurrence of order $37$ contributed by R.~H.~Hardin. It is
false. The entry publishes $30$ terms and links a b-file running to $n=210$, so the recurrence
can be tested at $163$ indices and it fails at every one; the first is $n=38$, where it exceeds
the true value by $393$. A transfer matrix on $1461$ states --- independently confirmed by
direct enumeration and reproducing all $200$ b-file terms --- gives the recurrence that does
hold. It has order $43$, its first $37$ coefficients are exactly the published ones, and the
annihilation test certifies it from $n=44$ onwards while returning \emph{no threshold at all}
for the published line. So the published line is the correct recurrence with its last six terms
dropped.
\end{abstract}

\noindent\small 2020 Mathematics Subject Classification. 05A15, 05B45, 15B36.\normalsize

\section{The sequence and the claim}

OEIS A196074 is ``Number of nX4 0..4 arrays with each element x equal to the number its
horizontal and vertical neighbors equal to 0,3,2,1,4 for x=0,1,2,3,4.''. It has offset $1$,
publishes $30$ terms beginning
\[
0,\;0,\;1,\;8,\;39,\;60,\;111,\;343,\;861,\;2115,\;5016,\;13734,\ \dots
\]
and links a b-file of $200$. The entry carries one formula line, an empirical recurrence of
order $37$:
\begin{multline}\label{eq:conj}
a(n)\;=\;a(n-2) + 4\,a(n-3) + 14\,a(n-4) + 23\,a(n-5)\\
{}- 2\,a(n-6) + 24\,a(n-7) + 44\,a(n-8) - 133\,a(n-9)\\
{}- 186\,a(n-10) - 317\,a(n-11) - 363\,a(n-12) + 372\,a(n-13)\\
{}+ 825\,a(n-14) - 50\,a(n-15) - 513\,a(n-16) + 173\,a(n-17)\\
{}+ 42\,a(n-18) + 441\,a(n-19) + 835\,a(n-20) + 563\,a(n-21)\\
{}+ 336\,a(n-22) - 1995\,a(n-23) - 4044\,a(n-24) - 1095\,a(n-25)\\
{}+ 4154\,a(n-26) + 1656\,a(n-27) - 395\,a(n-28) + 3391\,a(n-29)\\
{}- 709\,a(n-30) - 3210\,a(n-31) - 1608\,a(n-32) - 2382\,a(n-33)\\
{}- 47\,a(n-34) - 219\,a(n-35) + 91\,a(n-36) + 1032\,a(n-37).
\end{multline}
As of the live entry (revision $9$, Oct 2025) this is recorded as empirical, with no range
stated and nothing recording it as settled or corrected.

With offset $1$ and order $37$ the line first asserts something at $n=38$. DATA stops at
$n=30$, so on the DATA alone it is untestable --- but the b-file tests it $163$ times.

\section{The published recurrence fails at all 163 indices its b-file can test}

Substituting the b-file's own integers into the right-hand side of \eqref{eq:conj}:

\begin{center}\small
\begin{tabular}{rrrr}
\toprule
$n$ & $a(n)$ & \eqref{eq:conj} predicts & difference \\
\midrule
$38$ & $299231374661523$ & $299231374661916$ & $-393$\\
$39$ & $748157125706655$ & $748157125707069$ & $-414$\\
$40$ & $1870872919049913$ & $1870872919049599$ & $+314$\\
$41$ & $4678203678514716$ & $4678203678514596$ & $+120$\\
$42$ & $11697062923684722$ & $11697062923682370$ & $+2352$\\
$43$ & $29247168388430532$ & $29247168388418874$ & $+11658$\\
\bottomrule
\end{tabular}
\end{center}

\begin{proposition}\label{prop:false}
\eqref{eq:conj} is false. It fails at every one of the $163$ indices $n=38,\dots,200$ at
which the entry's b-file can test it. The smallest is $n=38$, where
$a(38)=299231374661523$ and \eqref{eq:conj} gives $299231374661916$, exceeding it by $393$.
\end{proposition}

\begin{proof}
Direct evaluation in exact integer arithmetic on the $200$ integers the entry's b-file
publishes. No model enters this computation.
\end{proof}

\section{What the entry counts, and a model for it}

Write $f=(0,3,2,1,4)$, that is $f(0)=0$, $f(1)=3$, $f(2)=2$, $f(3)=1$, $f(4)=4$. An array is
an $n\times4$ matrix over $\{0,1,2,3,4\}$, and the condition is that every cell of value $x$
has exactly $x$ of its horizontal and vertical neighbours equal to $f(x)$, positions outside
the array not being counted. The entry states this in the same words:
``Every 0 is next to 0 0's, every 1 is next to 1 3's, every 2 is next to 2 2's, every 3 is
next to 3 1's, every 4 is next to 4 4's''.

Because the neighbourhood reaches one row up and one row down, a state carrying a single row
is not enough. Take ordered pairs of consecutive rows as vertices, with $(p,c)\to(c,x)$ an
edge when every cell of the middle row $c$ satisfies the condition with $p$ above and $x$
below, and mark the pairs whose first row passes with nothing above and those whose second
passes with nothing below. Of the $5^{8}=390625$ ordered pairs only $S=1461$ are reachable,
and merging states with identical future behaviour leaves $58$.

The model reproduces all $30$ published DATA terms and all $200$ b-file terms exactly. It was
not taken on trust: an independent enumeration, written from the entry's wording alone and
sharing no code with the model, gives
\[
0,\;0,\;1,\;8,\;39,\;60,\;111 \qquad (n=1,\dots,7),
\]
the entry's first seven terms.

\section{The recurrence that does hold}

Being a walk count, $a$ is $C$-finite. Solving over $\mathbf{Q}$ from the b-file's own terms,
the least order that admits a recurrence is $43$, and the coefficients come out integral:
\begin{multline}\label{eq:true}
a(n)\;=\;a(n-2) + 4\,a(n-3) + 14\,a(n-4) + 23\,a(n-5)\\
{}- 2\,a(n-6) + 24\,a(n-7) + 44\,a(n-8) - 133\,a(n-9)\\
{}- 186\,a(n-10) - 317\,a(n-11) - 363\,a(n-12) + 372\,a(n-13)\\
{}+ 825\,a(n-14) - 50\,a(n-15) - 513\,a(n-16) + 173\,a(n-17)\\
{}+ 42\,a(n-18) + 441\,a(n-19) + 835\,a(n-20) + 563\,a(n-21)\\
{}+ 336\,a(n-22) - 1995\,a(n-23) - 4044\,a(n-24) - 1095\,a(n-25)\\
{}+ 4154\,a(n-26) + 1656\,a(n-27) - 395\,a(n-28) + 3391\,a(n-29)\\
{}- 709\,a(n-30) - 3210\,a(n-31) - 1608\,a(n-32) - 2382\,a(n-33)\\
{}- 47\,a(n-34) - 219\,a(n-35) + 91\,a(n-36) + 1032\,a(n-37)\\
{}274\,a(n-38) + 112\,a(n-39) + 84\,a(n-40) - 56\,a(n-41)\\
{}+ 8\,a(n-42) - 8\,a(n-43).
\end{multline}
The first $37$ coefficients of \eqref{eq:true} are, term for term, those of \eqref{eq:conj}.
The published line is \eqref{eq:true} with its last six terms,
$+274\,a(n-38)+112\,a(n-39)+84\,a(n-40)-56\,a(n-41)+8\,a(n-42)-8\,a(n-43)$, missing.

\begin{theorem}\label{thm:true}
\eqref{eq:true} holds for every $n>42$.
\end{theorem}

\begin{proof}
Let $q(t)=t^{43}-\sum_i c_it^{43-i}$ be its characteristic polynomial and $w=q(M)\mathbf u$,
with $M$ the transfer matrix above and $\mathbf u$ the start vector. For $n$ beyond the order,
\[
a(n)-\sum_i c_i\,a(n-i)\;=\;\mathbf v^{\!\top}M^{\,n-44}w ,
\]
so the recurrence holds from there on exactly when $\mathbf v^{\!\top}M^{m}w=0$ for every
$m\ge0$; by the Cayley--Hamilton theorem $M^{S}$ is an integer combination of
$I,M,\dots,M^{S-1}$, so checking $m<S$ suffices. Carried out in exact integer arithmetic on
the merged model, every value vanishes, and the test returns the threshold $42$.
\end{proof}

\begin{corollary}
\eqref{eq:conj} admits no threshold: there is no index beyond which it holds.
\end{corollary}

\begin{proof}
The same test applied to the published coefficients returns no threshold --- the residuals
$\mathbf v^{\!\top}M^{m}w$ for the order-$37$ polynomial do not all vanish, at any $m$ within
the Cayley--Hamilton bound, so they do not vanish eventually. Proposition~\ref{prop:false}
exhibits the failure at $163$ consecutive indices in any case.
\end{proof}

\section{Verification}

Five checks.

First, the disproof of Section~2 uses only the $200$ integers of the entry's own b-file, and
does not depend on the model or on Theorem~\ref{thm:true}.

Second, the model was confirmed against an independent enumeration written from the entry's
wording alone, agreeing at $n=1,\dots,7$, and it reproduces all $30$ DATA terms and all $200$
b-file terms --- every one, not a sample.

Third, \eqref{eq:true} was checked directly against the b-file at all $157$ indices from
$n=44$ to $n=200$ at which it asserts anything, and holds at every one.

Fourth, every order below $43$ was solved for and each fails, so $43$ is minimal and the
published $37$ is not merely incomplete bookkeeping but too short to be repaired by any
choice of coefficients.

Fifth, the annihilation test was run on both lines through the same code: it certifies
\eqref{eq:true} from $n=44$ and returns nothing for \eqref{eq:conj}.

\begin{remark}
This is the third entry of this shape settled here in one day, after A269637 (order 10
published, 13 true) and A236647 (34 published, 38 true). In all three the published
coefficients are not wrong but incomplete, and in all three they are the true recurrence's
first coefficients term for term. A recurrence fitted to a finite sample and reported at an
order the sample cannot support looks exactly like this: correct as far as it goes, and false
at the first index where it says anything.
\end{remark}

\begin{thebibliography}{9}
\bibitem{oeis} The OEIS Foundation, \emph{The On-Line Encyclopedia of Integer Sequences},
\url{https://oeis.org}, 2026. Sequence A196074.
\bibitem{stanley} R.~P.~Stanley, \emph{Enumerative Combinatorics}, Volume 1, 2nd ed.,
Cambridge University Press, 2011. (Transfer-matrix method, Section 4.7.)
\bibitem{kp} M.~Kauers and P.~Paule, \emph{The Concrete Tetrahedron}, Springer, 2011.
($C$-finite sequences and their recurrences.)
\bibitem{horn} R.~A.~Horn and C.~R.~Johnson, \emph{Matrix Analysis}, 2nd ed., Cambridge
University Press, 2013.
\end{thebibliography}

\end{document}
